\section{Typing}
\label{sec:b-typing}

Both typing developments share the same Lean type of B types:
\[
  \mathsf{BType} ::= \mathsf{int}
  \mid \mathsf{bool}
  \mid \mathsf{set}(\mathsf{BType})
  \mid \mathsf{prod}(\mathsf{BType},\mathsf{BType}).
\]
The conversion from a B type to the B expression denoting its carrier set is
also formalized:
\[
\begin{array}{rcl}
\mathsf{toTerm}(\mathsf{int}) &=& \mathbb{Z},\\
\mathsf{toTerm}(\mathsf{bool}) &=& \mathbb{B},\\
\mathsf{toTerm}(\mathsf{set}\,\tau) &=& \powB\mathsf{toTerm}(\tau),\\
\mathsf{toTerm}(\mathsf{prod}\,\tau\,\sigma) &=&
  \mathsf{toTerm}(\tau) \prodB \mathsf{toTerm}(\sigma).
\end{array}
\]
This function is used when a type must be reflected as a set-valued B
expression.

\subsection{Concrete Typing}
\label{subsec:b-concrete-typing}

Concrete typing is defined over the first-order syntax.
A typing context is a finite association list from variable names to B types,
and the judgment is the inductive predicate
\[
  \mathsf{Typing} :
  \mathsf{TypeContext} \to \mathsf{Term} \to \mathsf{BType} \to \mathsf{Prop},
\]
written in the Lean sources as \(\Gamma \vdash^B t : \tau\).

The rules are syntax-directed.
For each concrete constructor, the corresponding constructor of
\(\mathsf{Typing}\) stores exactly the premises needed to establish the result
type.
For example, \(\texttt{Typing.union}\) requires both operands to have the same
set type, while \(\texttt{Typing.app}\) requires the function expression to
have type \(\setB(\tau \prodB \sigma)\) and the argument to have type \(\tau\).
Binder rules carry the side conditions that are implicit in informal
presentations: the binder must bind at least one variable, the variable list
must not contain duplicates, the bound variables must be fresh for the
surrounding context, and the lists of variables, domains, and types must have
the same length.

Several structural facts are proved alongside the concrete rules.
The most important for later chapters is type determinism:
\[
  \btype{\Gamma}{t}{\tau}
  \quad\land\quad
  \btype{\Gamma}{t}{\sigma}
  \quad\Longrightarrow\quad
  \tau = \sigma.
\]
The first-order development proves this as \(\texttt{Typing.det}\).
The development also proves inversion lemmas for each constructor, such as
extracting the operand types from a typing derivation of a union, application,
lambda, or comprehension.
These lemmas turn a typing derivation into the precise type information needed
by later recursive arguments.

\subsection{Abstract Typing}
\label{subsec:b-abstract-typing}

Abstract typing is defined over the PHOAS syntax.
Its context is a function
\(\mathcal V \to \mathsf{Option}\,\mathsf{BType}\), where
\(\mathsf{none}\) represents absence from the context and
\(\mathsf{some}\,\tau\) represents a successful lookup.
The PHOAS judgment has the form
\[
  \mathsf{Typing} : (\mathcal V \to \mathsf{Option}\,\mathsf{BType})
  \to \mathsf{Term}\,\mathcal V \to \mathsf{BType} \to \mathsf{Prop},
\]
written \(\Gamma \vdash^{B'} t : \tau\).

The non-binding rules mirror the concrete typing rules.
The binder rules differ in how they expose the bound variables: instead of
checking a list of names and extending a finite association list, the rule
receives a family of fresh PHOAS variables and types the body under the
corresponding functional context extension.
Thus freshness and substitution bookkeeping are delegated to Lean's own
function space, while the B-specific typing premises remain explicit.

The PHOAS development contains the analogues of the concrete structural lemmas,
including type determinism under the expected nonemptiness assumption for
variable families.
These results are used by the semantics and by the SMT encoding proof, where
recursive calls require precise information about the type assigned to each
subterm.
